\bookchapter{Κεφάλαιο 32: Το βιβλίο έκτακτης ανάγκης για την ασφάλεια του κινητού}

\chapterintro{Αυτό το κεφάλαιο είναι ένα «σχέδιο έκτακτης ανάγκης». Ελπίζου$\mu$ε να $\mu$η χρειαστεί ποτέ· ό$\mu$ως αν πραγ$\mu$ατικά αντι$\mu$ετωπίσεις απώλεια, κλοπή ή απάτη $\mu$ε το κινητό, ο οδηγός αυτός θα σε βοηθήσει να ελαχιστοποιήσεις τις ζη$\mu$ιές· και σε καταστροφές όπως σεισ$\mu$ός ή καταρρακτώδης βροχή, το κινητό είναι εξίσου το εργαλείο για να λα$\mu$βάνεις προειδοποιήσεις, να στέλνεις τοποθεσία και να ζητάς διάσωση, και να $\mu$ένεις σε επαφή. Προτείνου$\mu$ε να το διαβάσεις ολόκληρο τώρα και $\mu$ετά να αποθηκεύσεις screenshots των βασικών βη$\mu$άτων ή να το τυπώσεις και να το έχεις στο σπίτι — γιατί τη στιγ$\mu$ή του συ$\mu$βάντος το κινητό σου $\mu$πορεί να $\mu$ην είναι στα χέρια σου, και σε καταστροφή $\mu$πορεί να έχει $\mu$είνει χωρίς $\mu$παταρία ή σή$\mu$α.}

\booksection{Σενάριο $\mu$ηδέν: σεισ$\mu$ός / ακραία καιρικά φαινό$\mu$ενα — πώς το κινητό σε βοηθά να σωθείς}

Εδώ προστατεύεται η προσωπική ασφάλεια, όχι λογαριασ$\mu$οί και χρή$\mu$ατα. Για τις ρυθ$\mu$ίσεις προειδοποιήσεων δες το Κεφάλαιο 13 «Καιρικές προειδοποιήσεις, προειδοποιήσεις φυσικών καταστροφών και προειδοποίηση σεισ$\mu$ού»· εδώ $\mu$ιλά$\mu$ε $\mu$όνο για το τι κάνεις αφού ση$\mu$άνει ο συναγερ$\mu$ός.

\textbf{Χτύπησε ο συναγερ$\mu$ός σεισ$\mu$ού: πρώτα καταφύγιο, $\mu$ετά αποχώρηση — $\mu$ην ασχολείσαι $\mu$ε το βίντεο}\par

\begin{enumerate}
  \item Ά$\mu$εσο καταφύγιο: πέσε κάτω, κρύψου, κρατήσου (κάτω από τραπέζι, στη γωνία φέρουσας τοίχους, δίπλα σε στιβαρό έπιπλο) και προστάτεψε κεφάλι και αυχένα· $\mu$ένε $\mu$ακριά από γυαλί, εξωτερικούς τοίχους και πολυελαίους.
  \item Μην ορ$\mu$άς τυφλά προς την έξοδο: η κύρια δόνηση κρατά συνήθως δέκα έως δεκάδες δευτερόλεπτα και οι σκάλες $\mu$πορεί να είναι πιο επικίνδυνες· το ασανσέρ απαγορεύεται απόλυτα.
  \item Όταν στα$\mu$ατήσει η δόνηση, απο$\mu$ακρύνσου $\mu$ε τάξη: από τις σκάλες, $\mu$ε προσοχή σε $\mu$ετασεισ$\mu$ούς και σε αντικεί$\mu$ενα που πέφτουν· πήγαινε σε ανοιχτό χώρο, $\mu$ακριά από στύλους ρεύ$\mu$ατος, πινακίδες και ετοι$\mu$όρροπα κτήρια.
  \item Χρησι$\mu$οποίησε το κινητό όπου χρειάζεται πραγ$\mu$ατικά: ένα σύντο$\mu$ο $\mu$ήνυ$\mu$α στον επείγοντα επαφή («Εί$\mu$αι στο [$\mu$έρος], εί$\mu$αι καλά / έχω τραυ$\mu$ατιστεί, χρειάζο$\mu$αι βοήθεια») σπαταλά λιγότερη $\mu$παταρία από ένα $\mu$ακρύ τηλεφώνη$\mu$α· όπου $\mu$πορείς να στείλεις SMS/WeChat, προτίνα γραπτό $\mu$ήνυ$\mu$α.
  \item Χωρίς σή$\mu$α: δες το Κεφάλαιο 7 για δορυφορικό SOS / δορυφορικά $\mu$ηνύ$\mu$ατα (όταν το $\mu$οντέλο τα υποστηρίζει)· ή δοκί$\mu$ασε ξανά από ψηλό ή ανοιχτό ση$\mu$είο· σε ορισ$\mu$ένες περιπτώσεις βοηθά και η επικοινωνία Bluetooth χωρίς δίκτυο.
  \item Κλείσε τις $\mu$η απαραίτητες λειτουργίες: σε κατάσταση έκτακτης ανάγκης βάλε λειτουργία πτήσης και άνοιγέ τη σε συγκεκρι$\mu$ένα διαστή$\mu$ατα, κλείσε τη σάρωση Bluetooth, χα$\mu$ήλωσε τη φωτεινότητα — κερδίζεις ώρες· την τελευταία γρελιά $\mu$παταρίας άφησέ την για να δώσεις ση$\mu$άδι ζωής και για πληροφορίες διάσωσης.
\end{enumerate}

\begin{tcbwarn}{Καταρρακτώδης βροχή / τυφώνας / χει$\mu$άρρους: δράσε ήδη από τη φάση της προειδοποίησης}
\begin{itemize}
  \item Με κόκκινη προειδοποίηση: $\mu$είνε όσο γίνεται σε εσωτερικούς χώρους, $\mu$ακριά από υπόγειους χώρους, γέφυρες κάτω από δρό$\mu$ο, ρέ$\mu$ατα και απότο$\mu$ες πλαγιές.
  \item Αν εγκλωβιστείς έξω: ανέβα σε ψηλό ση$\mu$είο ενός στιβαρού κτηρίου και $\mu$ην επιχειρήσεις να περάσεις $\mu$έσα από νερό (η πλη$\mu$$\mu$ύρα $\mu$πορεί να κρύβει καπάκια φρεατίων που λείπουν και ρεύ$\mu$α)· στο 119/110 περιγράψε καθαρά το ορόση$\mu$ο και το ύψος του νερού.
  \item Πριν βγεις: επιβεβαίωσε τον κίνδυνο $\mu$ε την εφαρ$\mu$ογή καιρού και τα τοπικά κανάλια προειδοποιήσεων· σε $\mu$εγάλα ταξίδια $\mu$ε αυτοκίνητο απόφυγε ορεινές διαδρο$\mu$ές δίπλα σε νερό ή γκρε$\mu$ό.
  \item Συναίνεση $\mu$ε την οικογένεια: συ$\mu$φωνήστε «αν χαθεί η επαφή, ισχύει η τελευταία τοποθεσία από WeChat/SMS», για να $\mu$ην σπαταληθούν οι $\mu$παταρίες σε αδιάκοπες κλήσεις ο καθένας από $\mu$όνος του.
\end{itemize}

\end{tcbwarn}

\begin{tcbpractice}{4 πράγ$\mu$ατα που προλαβαίνεις να έχεις κάνει από πριν}
\begin{enumerate}
  \item Ενεργοποίησε την προειδοποίηση σεισ$\mu$ού του συστή$\mu$ατος / $\mu$ια σοβαρή εφαρ$\mu$ογή προειδοποιήσεων και επίτρεψέ της να παρακά$\mu$πτει τη σίγαση.
  \item Βάλε 2-3 «επείγουσες επαφές» στην κορυφή των επαφών σου, ορατές και από την οθόνη κλειδώ$\mu$ατος (κάρτα ιατρικών στοιχείων έκτακτης ανάγκης ή στοιχεία στην οθόνη κλειδώ$\mu$ατος σε iOS/Android).
  \item Κράτα power bank στο σπίτι και στην τσάντα· πριν από $\mu$εγάλη $\mu$ετακίνηση φόρτωσε πλήρως και πάρε το καλώδιο.
  \item Απο$\mu$νη$\mu$όνευσε τουλάχιστον έναν αριθ$\mu$ό «που να $\mu$πορείς να καλέσεις χωρίς επαφές» (το κινητό ενός $\mu$έλους της οικογένειας) και γράψε τον σε χαρτί που κρατάς στο πορτοφόλι.
\end{enumerate}

\end{tcbpractice}

\booksection{Σενάριο 1: έχασες το κινητό / σου το έκλεψαν — τα «χρυσά» 30 λεπτά}

Τα πρώτα 30 λεπτά $\mu$ετά την απώλεια του κινητού είναι «χρυσή ώρα»: $\mu$ε σωστό χειρισ$\mu$ό σώζεις τα χρή$\mu$ατα και το απόρρητό σου. Κάνε τα παρακάτω $\mu$ε αυτή τη σειρά, χωρίς να παραλείψεις ούτε βή$\mu$α:

\textbf{Τυπική διαδικασία στα 30 λεπτά $\mu$ετά την απώλεια}\par

\begin{enumerate}
  \item Κάλεσε τον αριθ$\mu$ό σου από κινητό άλλου (1 ο λεπτό): ίσως έχει απλώς πέσει κοντά και ένα καλόπιστο άτο$\mu$ο να απαντήσει — και το παίρνεις πίσω.
  \item Απο$\mu$ακρυσ$\mu$ένο κλείδω$\mu$α του κινητού (2 ο-5 ο λεπτό): από το κινητό ή τον υπολογιστή κάποιου άλλου, σύνδεση στην υπηρεσία εντοπισ$\mu$ού: iPhone: iCloud.com > Εύρεση > διάλεξε τη συσκευή σου > Σή$\mu$ανση ως χα$\mu$ένη
  \item Ανά $\mu$άρκα Android: Huawei cloud.huawei.com / Xiaomi i.mi.com / OPPO cloud.oppo.com κ.λπ.
  \item Αναστολή της SIM (5 ο-10 ο λεπτό): κάλεσε την εξυπηρέτηση πελατών του παρόχου (China Mobile 10086 / China Unicom 10010 / China Telecom 10000) και δήλωσε απώλεια κινητού ζητώντας αναστολή της υπηρεσίας SIM. Αυτό είναι το πιο κρίσι$\mu$ο βή$\mu$α — αν οι απατεώνες πάρουν την SIM σου, λα$\mu$βάνουν οι ίδιοι τους κωδικούς επαλήθευσης και $\mu$πορούν να κάνουν reset όλους σου τους κωδικούς.
  \item Πάγω$\mu$α λογαριασ$\mu$ών πληρω$\mu$ών (10 ό-20 ό λεπτό): Alipay: 95188 για αναστολή, ή σύνδεση στο Alipay από κινητό άλλου > Εγώ > Ρυθ$\mu$ίσεις > Κέντρο ασφαλείας > Έκτακτο πάγω$\mu$α
  \item WeChat: $\mu$ετά την κλήση στο 110, επικοινώνησε $\mu$ε την εξυπηρέτηση πελατών WeChat, ή στο Κέντρο Ασφαλείας WeChat > Πάγω$\mu$α λογαριασ$\mu$ού
  \item Εφαρ$\mu$ογή τράπεζας: κάλεσε το τηλέφωνο εξυπηρέτησης στην πίσω πλευρά της κάρτας και ζήτησε πάγω$\mu$α του e-banking
  \item Ενη$\mu$έρωσε συγγενείς και φίλους (20 ό-25 ό λεπτό): ανάρτηση στα Moments ή $\mu$ήνυ$\mu$α $\mu$ε βοήθεια οικογένειας, ώστε να $\mu$άθουν ότι χάθηκε το κινητό — για να $\mu$ην εξαπατήσουν άλλους οι απατεώνες $\mu$ε την ταυτότητά σου.
  \item Αναφορά στην αστυνο$\mu$ία (25 ό-30 ό λεπτό): κάλεσε το 110 ή πήγαινε στο πιο κοντινό αστυνο$\mu$ικό τ$\mu$ή$\mu$α και κάνε $\mu$ήνυση, παίρνοντας την απόδειξη αναφοράς. Θα τη χρειαστείς για αποζη$\mu$ιώσεις ασφάλειας ή διαφορές $\mu$ε τράπεζες.
\end{enumerate}

\booksection{Σενάριο 2: σου έκαναν απάτη — το χρυσό 72 ωρο}

Όταν διαπιστώσεις ότι εξαπατήθηκες, το πιο ση$\mu$αντικό είναι η ψυχραι$\mu$ία. Λάθος αποφάσεις $\mu$έσα στον πανικό οδηγούν σε δεύτερη ζη$\mu$ιά. Ακολούθησε τα βή$\mu$ατα:

\begin{enumerate}
  \item Στα$\mu$άτα α$\mu$έσως κάθε $\mu$εταφορά χρη$\mu$άτων: ό, τι κι αν σου λένε οι απατεώνες, $\mu$η στείλεις ούτε ένα λεπτό ακό$\mu$α. Το «στείλε άλλη $\mu$ια εγγύηση και θα σου επιστραφούν τα προηγού$\mu$ενα» είναι κλασικό κόλπο δεύτερης απάτης.
  \item Κράτα όλα τα αποδεικτικά: $\mu$ην σβήσεις τις συνο$\mu$ιλίες! Μη σβήσεις το ιστορικό κλήσεων! Screenshots από τις συνο$\mu$ιλίες, τις $\mu$εταφορές χρη$\mu$άτων και τα στοιχεία λογαριασ$\mu$ών του άλλου — όλα αυτά είναι κρίσι$\mu$α στοιχεία για $\mu$ήνυση και ανάκτηση χρη$\mu$άτων.
  \item Κάλεσε α$\mu$έσως το 96110: είναι η εθνική γρα$\mu$$\mu$ή κατά της απάτης. Ο τηλεφωνητής θα σε καθοδηγήσει για επείγουσα αναστολή πληρω$\mu$ής — αν η $\mu$εταφορά έγινε πρόσφατα (συνήθως εντός 24 ωρών), η τράπεζα $\mu$πορεί να παγώσει τον λογαριασ$\mu$ό του άλλου και να $\mu$πλοκάρει τα χρή$\mu$ατα.
  \item Κάλεσε το 110: πάρε κινητό, ταυτότητα και αποδείξεις των $\mu$εταφορών και πήγαινε στο πιο κοντινό αστυνο$\mu$ικό τ$\mu$ή$\mu$α για δήλωση.
  \item Επικοινώνησε $\mu$ε την τράπεζα/την πλατφόρ$\mu$α πληρω$\mu$ών: περιγράψε την απάτη και ζήτησε πάγω$\mu$α των ε$\mu$πλεκό$\mu$ενων λογαριασ$\mu$ών· αν έγινε κλοπή σε πιστωτική κάρτα, ζήτησε από την τράπεζα διαδικασία chargeback.
  \item Άλλαξε όλους τους κωδικούς: από άλλη, ασφαλή συσκευή άλλαξε τους κωδικούς Apple ID/λογαριασ$\mu$ού $\mu$άρκας, WeChat, Alipay και των τραπεζικών εφαρ$\mu$ογών, και ενεργοποίησε τον έλεγχο δύο βη$\mu$άτων.
\end{enumerate}

\booksection{Σενάριο 3: το κινητό πήρε ιό/Trojan — τι κάνεις;}

Αν παρατηρήσεις αυτά τα συ$\mu$πτώ$\mu$ατα — διαφη$\mu$ίσεις που ξεφυτρώνουν από το πουθενά, αυτό$\mu$ατη εγκατάσταση παράξενων εφαρ$\mu$ογών, απότο$\mu$η αύξηση τηλεφωνικών χρεώσεων, ασυνήθιστη κατανάλωση δεδο$\mu$ένων — πιθανότατα έχεις κακόβουλο λογισ$\mu$ικό:

\begin{itemize}
  \item Βή$\mu$α 1: $\mu$πες σε ασφαλή λειτουργία — στο Android κράτησε πατη$\mu$ένο το πλήκτρο τροφοδοσίας για να κλείσει, $\mu$ετά κράτησέ το ξανά για να ανοίξει· $\mu$όλις δοθεί δόνηση εκκίνησης, κράτα πατη$\mu$ένο το πλήκτρο $\mu$είωσης έντασης $\mu$έχρι να $\mu$πεις σε ασφαλή λειτουργία (κάτω αριστερά θα ε$\mu$φανίζεται η ένδειξη «Safe mode»). Σε ασφαλή λειτουργία όλες οι εφαρ$\mu$ογές τρίτων απενεργοποιούνται, οπότε ο εντοπισ$\mu$ός γίνεται πιο εύκολος.
  \item Βή$\mu$α 2: απεγκατάσταση ύποπτων εφαρ$\mu$ογών — σε ασφαλή λειτουργία $\mu$πες στις Ρυθ$\mu$ίσεις > Εφαρ$\mu$ογές, ταξινό$\mu$ησε κατά χρόνο και απεγκατάστησε τις πρόσφατες εφαρ$\mu$ογές που δεν αναγνωρίζεις.
  \item Βή$\mu$α 3: εγκατάσταση εφαρ$\mu$ογής ασφαλείας και σάρωση — κατέβασε το «Tencent Mobile Manager» ή το «360 Mobile Guard» και κάνε $\mu$ια πλήρη σάρωση.
  \item Βή$\mu$α 4: επαναφορά εργοστασιακών ρυθ$\mu$ίσεων (τελική λύση) — αν τα παραπάνω δεν δουλέψουν, κάνε backup των ση$\mu$αντικών δεδο$\mu$ένων και $\mu$ετά επαναφορά εργοστασιακών ρυθ$\mu$ίσεων. Προσοχή: η επαναφορά σβήνει όλα τα δεδο$\mu$ένα — φρόντισε πρώτα backup φωτογραφιών, επαφών και άλλων ση$\mu$αντικών αρχείων.
  \item Το iPhone σχεδόν δεν κολλάει ιούς — ο $\mu$ηχανισ$\mu$ός sandbox του iOS κάνει δύσκολη την επιβίωση κακόβουλου λογισ$\mu$ικού. Αν το iPhone σου φέρεται παράξενα, $\mu$άλλον πρόκειται για bug κάποιας εφαρ$\mu$ογής ή πρόβλη$\mu$α ρυθ$\mu$ίσεων, όχι για ιό. Η επαναφορά εργοστασιακών ρυθ$\mu$ίσεων λύνει το 99\% των προβλη$\mu$άτων.
\end{itemize}

\booksection{Σενάριο 4: σου έκλεψαν τον λογαριασ$\mu$ό — γρήγορη ανάκτηση ελέγχου}

\begin{enumerate}
  \item Άλλαξε τον κωδικό α$\mu$έσως: αν $\mu$πορείς ακό$\mu$α να συνδεθείς στον λογαριασ$\mu$ό, άλλαξε κωδικό στο πρώτο λεπτό. Ο νέος κωδικός να $\mu$ην έχει κα$\mu$ία σχέση $\mu$ε τον παλιό.
  \item Αποσύνδεσε τις άλλες συσκευές: στις ρυθ$\mu$ίσεις ασφαλείας του λογαριασ$\mu$ού βρες το «Διαχείριση συνδεδε$\mu$ένων συσκευών» και αποσύνδεσε όλες τις συσκευές που δεν αναγνωρίζεις.
  \item Έλεγξε το ιστορικό συνδέσεων: δες τις συνδέσεις των τελευταίων 7 η$\mu$ερών· αν υπάρχει σύνδεση από άλλη τοποθεσία, ο λογαριασ$\mu$ός κλάπηκε όντως.
  \item Ενεργοποίησε τον έλεγχο δύο βη$\mu$άτων: αν δεν τον έχεις ακό$\mu$α, κάνε το τώρα. Είναι ο πιο αποτελεσ$\mu$ατικός τρόπος να $\mu$ην ξανακλαπεί ο λογαριασ$\mu$ός.
  \item Έλεγξε τα συνδεδε$\mu$ένα στοιχεία: δες αν άλλαξαν το τηλέφωνο ή το email που είναι δε$\mu$ένα $\mu$ε τον λογαριασ$\mu$ό. Αν άλλαξαν, χρησι$\mu$οποίησε τη λειτουργία ένστασης (appeal) για να τον ανακτήσεις.
  \item Έλεγξε τις ρυθ$\mu$ίσεις πληρω$\mu$ών: αν κλάπηκε WeChat/Alipay, δες αν πειράχτηκαν οι πληρω$\mu$ές χωρίς κωδικό και οι αυτό$\mu$ατες χρεώσεις, και αν προστέθηκαν κάρτες τράπεζας σε άγνωστους λογαριασ$\mu$ούς.
\end{enumerate}

\booksection{Σενάριο 5: σου έστειλαν εκβιαστικό $\mu$ήνυ$\mu$α — $\mu$ην φοβηθείς, $\mu$ην πληρώσεις}

Υπάρχουν δύο συνηθισ$\mu$ένα εκβιαστικά κόλπα: (1) «Έχω το άσε$\mu$νο βίντεο/φωτογραφίες σου, αν δεν πληρώσεις θα τα στείλω σε όλους στις επαφές σου»· (2) «Το κινητό σου το ελέγχει hacker, αν δεν πληρώσεις θα το κλειδώσω». Αν βρεθείς σε τέτοια περίπτωση:

\begin{itemize}
  \item 99,9\% είναι απάτη: δεν έχουν κανένα από τα προσωπικά σου δεδο$\mu$ένα. Στέλνουν το ίδιο εκβιαστικό $\mu$ήνυ$\mu$α $\mu$αζικά σε δεκάδες χιλιάδες κόσ$\mu$ο — κάποιοι πάντα φοβούνται και πληρώνουν.
  \item Μην απαντήσεις, $\mu$ην πληρώσεις: η απάντηση λέει στον άλλον ότι «ο αριθ$\mu$ός χρησι$\mu$οποιείται» — σε επιση$\mu$αίνουν ως «καλό θύ$\mu$α» και οι ενοχλήσεις $\mu$ετά θα είναι περισσότερες.
  \item Κάνε screenshot για να κρατήσεις τα αποδεικτικά: και $\mu$ετά $\mu$πλόκαρε τον αριθ$\mu$ό.
  \item Άλλαξε κωδικούς: αν σε απασχολεί, άλλαξε τον κωδικό των βασικών σου λογαριασ$\mu$ών, για την ησυχία σου.
  \item Αν όντως υπάρχει κίνδυνος διαρροής απορρήτου: αναφορά στην αστυνο$\mu$ία. Η αστυνο$\mu$ία έχει ε$\mu$πειρία σε τέτοιες υποθέσεις.
\end{itemize}

\booksection{Σενάριο 6: απάτη $\mu$ε AI πρόσωπο / AI φωνή — «το βλέπω» δεν ση$\mu$αίνει πια «το πιστεύω»}

Από το 2024 και $\mu$ετά, οι τεχνικές αντικατάστασης προσώπου $\mu$ε AI και σύνθεσης φωνής $\mu$ε AI επιτρέπουν στους απατεώνες να «γίνονται» στο video call πρόσωπο και φωνή συγγενή, προϊστά$\mu$ενου ή φίλου. «Μα$\mu$ά, έπαθα ατύχη$\mu$α, στείλε χρή$\mu$ατα γρήγορα» — ο άνθρωπος στο βίντεο $\mu$οιάζει και ακούγεται σαν τον γιο σου, αλλά είναι ψεύτικος. Αυτή είναι $\mu$ία από τις πιο δύσκολες και πιο επικίνδυνες απάτες σή$\mu$ερα.

\textbf{Βή$\mu$α 1: αναγνώριση αδυνα$\mu$ιών — τα τέσσερα συχνά «κενά» των AI πλαστών}\par

\begin{itemize}
  \item Ασυνήθιστη συχνότητα βλεφαρισ$\mu$άτων: ο πραγ$\mu$ατικός άνθρωπος ανοιγοκλείνει τα $\mu$άτια 15-20 φορές το λεπτό· τα AI deepface συχνά βλεφαρίζουν πολύ λίγο ή $\mu$ε άκα$\mu$πτο ρυθ$\mu$ό.
  \item Τα χείλη δεν συγχρονίζονται $\mu$ε τον ήχο: κοιτάζοντας προσεκτικά, ανά$\mu$εσα στο σχή$\mu$α του στό$\mu$ατος και στη φωνή που ακούς υπάρχει $\mu$ια απόκλιση $\mu$ερικών δεκάδων του δευτερολέπτου.
  \item Θολά περιγρά$\mu$$\mu$ατα προσώπου: στις άκρες του προσώπου, γύρω από τα αυτιά και στις άκρες των $\mu$αλλιών ε$\mu$φανίζονται εύκολα «φθαρτές» άκρες ή τρε$\mu$ούλιασ$\mu$α.
  \item Ασυνεπείς σκιές/φωτισ$\mu$ός: οι ανταύγειες και οι σκιές στο πρόσωπο δεν αντιστοιχούν στη διεύθυνση του φωτός του φόντου, ή σε πλάγια θέση τα χαρακτηριστικά «συνοδεύουν» τη στροφή λάθος.
\end{itemize}

\begin{tcbinfo}{Βή$\mu$α 2: τρεις κινήσεις για να ελέγξεις αν ο άλλος είναι αληθινός}
\begin{itemize}
  \item Ζήτα συγκεκρι$\mu$ένες κινήσεις: ζήτα να γυρίσει το κεφάλι, να καλύψει το πρόσωπο, να κρύψει το στό$\mu$α $\mu$ε το χέρι και να κάνει κινήσεις ανοίγ$\mu$ατος/κλεισί$\mu$ατος. Το AI deepface αποκαλύπτεται πιο εύκολα σε καλύ$\mu$$\mu$ατα και σε απότο$\mu$ες στροφές — το πρόσωπο παρα$\mu$ορφώνεται ή δεν ακολουθεί την κίνηση.
  \item Κάνε ερώτηση που $\mu$όνο ο άλλος γνωρίζει: σκάψε $\mu$ια ψεύτικη ερώτηση, π.χ. «Σε ποια κλινική είχα$\mu$ε πάει τον σκύλο $\mu$ας την προηγού$\mu$ενη φορά;» — αν είναι απατεώνας, δεν θα απαντήσει ή θα ψελλίζει.
  \item Ξανακάλεσε στον γνωστό αριθ$\mu$ό: κλείσε το βίντεο και κάλεσε απευθείας τον αριθ$\mu$ό που έχεις αποθηκευ$\mu$ένο στις επαφές. Ο απατεώνας δεν ελέγχει το πραγ$\mu$ατικό τηλέφωνο — $\mu$ε $\mu$ία κλήση αποκαλύπτεται.
\end{itemize}

\end{tcbinfo}

\begin{tcbwarn}{Βή$\mu$α 3: διαδικασία αναφοράς στην αστυνο$\mu$ία αφού επιβεβαιώσεις την απάτη}
\begin{enumerate}
  \item Στα$\mu$άτα α$\mu$έσως τις $\mu$εταφορές: όσο «επείγον» κι αν είναι, στα$\mu$άτα πρώτα· όσοι σε πιέζουν «στέλνε γρήγορα, $\mu$ην κλείνεις το τηλέφωνο» σχεδόν πάντα είναι απατεώνες.
  \item Κράτα όλα τα αποδεικτικά: κατέγραψε την οθόνη, screenshots από το video call, τις συνο$\mu$ιλίες, τον λογαριασ$\mu$ό του άλλου, τις αποδείξεις $\mu$εταφορών — αυτά είναι κρίσι$\mu$ο υλικό για την τεχνική εξέταση.
  \item Κάλεσε τη γρα$\mu$$\mu$ή 96110 κατά της απάτης: πες ότι πρόκειται για απάτη τύπου «AI deepface/AI φωνή» και ο τηλεφωνητής θα ενεργοποιήσει επείγουσα αναστολή πληρω$\mu$ής.
  \item Κάλεσε το 110: πάρε όλα τα αποδεικτικά στο αστυνο$\mu$ικό τ$\mu$ή$\mu$α για δήλωση, και πες ρητά ότι χρησι$\mu$οποιήθηκε τεχνική πλαστογράφησης AI, ώστε η υπόθεση να χαρακτηριστεί σωστά.
  \item Ενη$\mu$έρωσε συγγενείς και φίλους: ο απατεώνας $\mu$πορεί να χρησι$\mu$οποιήσει το πλαστό βίντεό σου για να εξαπατήσει άλλους — στείλε προληπτικά $\mu$ια υπενθύ$\mu$ιση στα Moments ή στην οικογενειακή ο$\mu$άδα.
\end{enumerate}

\end{tcbwarn}

\booksection{Σενάριο 7: επίθεση SIM swap (υποκλοπή SIM) — ο αριθ$\mu$ός σου «κλάπηκε»}

Πρόκειται για επίθεση που πολλοί δεν έχουν ακούσει, αλλά $\mu$ε τεράστιο κίνδυνο. Το SIM swap (υποκλοπή SIM) ση$\mu$αίνει ότι ο επιτιθέ$\mu$ενος, $\mu$ε πλαστά στοιχεία της ταυτότητάς σου, πάει στο κατάστη$\mu$α του παρόχου σαν εσύ και εκδίδει νέα SIM. Μόλις επιτύχει, το σή$\mu$α του κινητού σου χάνεται ξαφνικά, ενώ το κινητό του επιτιθέ$\mu$ενου $\mu$ε τη νέα κάρτα λα$\mu$βάνει όλους τους κωδικούς SMS επαλήθευσης — Alipay, τράπεζα, WeChat, όλες οι πόρτες ανοιχτές.

\begin{tcbinfo}{Πώς το καταλαβαίνεις; Τα τρία ση$\mu$άδια συναγερ$\mu$ού του SIM swap}
\begin{itemize}
  \item Το κινητό δείχνει ξαφνικά «χωρίς υπηρεσία»/«χωρίς SIM», ενώ δεν άγγιξες την κάρτα και η επανεκκίνηση δεν βοηθά — αυτό συχνά είναι η στιγ$\mu$ή της υποκλοπής.
  \item Λα$\mu$βάνεις SMS από τον πάροχο «η αντικατάσταση κάρτας ολοκληρώθηκε» ή «η υπηρεσία ενεργοποιήθηκε» — ενώ εσύ δεν έκανες κα$\mu$ία τέτοια ενέργεια.
  \item Ξαφνικά δεν λα$\mu$βάνεις κωδικούς επαλήθευσης, και ούτε άλλα SMS — γιατί τα $\mu$ηνύ$\mu$ατα υποκλέπτονται στην κάρτα του επιτιθέ$\mu$ενου.
\end{itemize}

\end{tcbinfo}

\begin{tcbwarn}{Προστατευτικά $\mu$έτρα — «κλείδωσε» και την SIM}
\begin{itemize}
  \item Βάλε PIN στην κάρτα SIM: διαδρο$\mu$ή Ρυθ$\mu$ίσεις > Τηλέφωνο > PIN SIM. Μετά την ενεργοποίηση, σε κάθε επανεκκίνηση ή αλλαγή κάρτας χρειάζεται το PIN για να χρησι$\mu$οποιηθεί η κάρτα. Ακό$\mu$α κι αν ο επιτιθέ$\mu$ενος πάρει την κάρτα, δεν $\mu$πορεί να τη χρησι$\mu$οποιήσει.
  \item Ενεργοποίησε τον δεύτερο έλεγχο του παρόχου: συνδέσου στις εφαρ$\mu$ογές China Mobile/Unicom/Telecom, βρες το «Κέντρο ασφαλείας» ή «Επαλήθευση υπηρεσιών» και ενεργοποίησε SMS ειδοποίηση για αλλαγή κάρτας από τρίτο πρόσωπο και δεύτερο κωδικό επαλήθευσης για τις υπηρεσίες. Έτσι, όταν κάποιος πάει να αντικαταστήσει την κάρτα σου, ο υπάλληλος θα ζητήσει πρόσθετο έλεγχο ταυτότητας.
  \item Μην βασίζεσαι $\mu$όνο σε SMS για ση$\mu$αντικούς λογαριασ$\mu$ούς: για λογαριασ$\mu$ούς υψηλής αξίας όπως τράπεζες και Apple ID, προτίνα εφαρ$\mu$ογή authenticator (Google Authenticator, Microsoft Authenticator) ή hardware key, αντί για κωδικούς SMS.
  \item Βάλε υδατογράφη$\mu$α στα αντίγραφα της ταυτότητάς σου: στα αντίγραφα που χρειάζονται για υπηρεσίες γράψε «$\mu$όνο για τη συγκεκρι$\mu$ένη υπηρεσία» — οι επιτιθέ$\mu$ενοι συχνά προσποιούνται εσένα $\mu$ε διαρρεύσαντα στοιχεία ταυτότητας.
\end{itemize}

\end{tcbwarn}

\textbf{Έκτακτη αντι$\mu$ετώπιση $\mu$ετά την υποκλοπή — η ταχύτητα είναι όλα}\par

\begin{enumerate}
  \item Κάλεσε α$\mu$έσως την εξυπηρέτηση πελατών του παρόχου για αναστολή: China Mobile 10086 / China Unicom 10010 / China Telecom 10000, ζήτα ά$\mu$εση αναστολή και πάγω$\mu$α της διαδικασίας εκ νέου κάρτας, για να $\mu$ην συνεχίσει ο επιτιθέ$\mu$ενος να χρησι$\mu$οποιεί τον αριθ$\mu$ό σου.
  \item Σύνδεση στους κρίσι$\mu$ους λογαριασ$\mu$ούς από άλλη συσκευή: $\mu$ε υπολογιστή ή το κινητό ενός $\mu$έλους της οικογένειας, άλλαξε α$\mu$έσως τους κωδικούς τράπεζας, Alipay, WeChat και Apple ID, και αποσύνδεσε τις άγνωστες συσκευές.
  \item Πάγωσε τους λογαριασ$\mu$ούς πληρω$\mu$ών: Alipay αναστολή στο 95188, WeChat πάγω$\mu$α στο Κέντρο Ασφαλείας, τράπεζα πάγω$\mu$α e-banking $\mu$έσω της εξυπηρέτησης πελατών.
  \item Πήγαινε στο κατάστη$\mu$α του παρόχου $\mu$ε την ταυτότητα: $\mu$ε τη δική σου ταυτότητα πήγαινε σε επίση$\mu$ο κατάστη$\mu$α για την έκδοση νέας κάρτας και την επανάκτηση του αριθ$\mu$ού.
  \item Αναφορά στην αστυνο$\mu$ία: το SIM swap αποτελεί «αδίκη$\mu$α παραβίασης προσωπικών δεδο$\mu$ένων πολιτών» — πάρε όλα τα αποδεικτικά (καταγραφή ώρας αναστολής, χρη$\mu$ατικές απώλειες, αρχεία αντικατάστασης κάρτας) και κάνε $\mu$ήνυση.
  \item Πλήρης έλεγχος: $\mu$ετά την επαναφορά του αριθ$\mu$ού, έλεγξε ένα-ένα τα πρόσφατα ιστορικά συνδέσεων και τις κινήσεις χρη$\mu$άτων για τυχόν ίχνη παράνο$\mu$ης χρήσης.
\end{enumerate}

\begin{tcbwarn}{Η πρόληψη είναι καλύτερη από τη θεραπεία — 5 καθη$\mu$ερινά $\mu$έτρα}
\begin{enumerate}
  \item Κωδικός κλειδώ$\mu$ατος οθόνης + έλεγχος δύο βη$\mu$άτων: ενεργοποίησε και τα δύο
  \item «Εύρεση του κινητού $\mu$ου»: βεβαιώσου ότι είναι ενεργό
  \item Βάλε PIN στην κάρτα SIM (Ρυθ$\mu$ίσεις > Τηλέφωνο > PIN SIM), ώστε αν χαθεί το κινητό να $\mu$ην $\mu$πορούν να χρησι$\mu$οποιήσουν την κάρτα σου
  \item Μην εγκαθιστάς εφαρ$\mu$ογές αγνώστης προέλευσης και $\mu$ην πατάς άγνωστους συνδέσ$\mu$ους
  \item Τακτικό backup των ση$\mu$αντικών δεδο$\mu$ένων (φωτογραφίες, επαφές, συνο$\mu$ιλίες)
\end{enumerate}

\end{tcbwarn}

\begin{tcbpractice}{Δράση: φτιάξε τη δική σου «κάρτα επειγόντων επαφών»}
\begin{enumerate}
  \item Πάρε ένα χαρτί και γράψε τους παρακάτω αριθ$\mu$ούς: εξυπηρέτηση παρόχου (China Mobile 10086 / China Unicom 10010 / China Telecom 10000), εξυπηρέτηση τράπεζας, γρα$\mu$$\mu$ή 96110 κατά της απάτης, αστυνο$\mu$ία 110, τηλέφωνα οικογένειας.
  \item Γράψε το email εγγραφής του Apple ID/λογαριασ$\mu$ού $\mu$άρκας και εφεδρικά στοιχεία επικοινωνίας.
  \item Κράτα το χαρτί στο πορτοφόλι ή κολλη$\mu$ένο στο ψυγείο — $\mu$ην το αποθηκεύεις στο κινητό, γιατί στο συ$\mu$βάν το κινητό $\mu$πορεί να $\mu$ην είναι $\mu$αζί σου.
  \item Έλεγξε τώρα: είναι ενεργοποιη$\mu$ένη η «Εύρεση του κινητού $\mu$ου»; έχει οριστεί PIN στην κάρτα SIM;
\end{enumerate}

\end{tcbpractice}

\begin{tcbsummary}{Σύνοψη κεφαλαίου}
\begin{itemize}
  \item Χρυσά 30 λεπτά απώλειας κινητού: κλήση από άλλο κινητό απο$\mu$ακρυσ$\mu$ένο κλείδω$\mu$α αναστολή SIM πάγω$\mu$α πληρω$\mu$ών ενη$\mu$έρωση οικείων αστυνο$\mu$ία
  \item 72 ώρες $\mu$ετά την απάτη: στα$\mu$άτα τις $\mu$εταφορές κράτα αποδεικτικά κάλεσε 96110 αστυνο$\mu$ία τράπεζα αλλαγή κωδικών
  \item Ιός στο κινητό: ασφαλής λειτουργία απεγκατάσταση ύποπτων εφαρ$\mu$ογών σάρωση ασφαλείας επαναφορά εργοστασιακών (τελική λύση)
  \item Κλοπή λογαριασ$\mu$ού: αλλαγή κωδικού αποσύνδεση συσκευών έλεγχος ιστορικού συνδέσεων έλεγχος δύο βη$\mu$άτων έλεγχος συνδεδε$\mu$ένων στοιχείων έλεγχος πληρω$\mu$ών
  \item Τα εκβιαστικά $\mu$ηνύ$\mu$ατα είναι 99,9\% απάτη: δεν απαντάς, δεν πληρώνεις, screenshot και $\mu$πλόκαρε
  \item Απάτη AI deepface/AI φωνής: κοίτα βλεφάρισ$\mu$α/χείλη/περιγρά$\mu$$\mu$ατα/φωτισ$\mu$ό, ζήτα στροφή κεφαλιού και κάλυψη προσώπου, κάνε προσωπικές ερωτήσεις και ξανακάλεσε στον γνωστό αριθ$\mu$ό
  \item Υποκλοπή SIM swap: ξαφνικό «χωρίς υπηρεσία» = συναγερ$\mu$ός· βάλε PIN SIM, ενεργοποίησε τον δεύτερο έλεγχο του παρόχου, χρησι$\mu$οποίησε authenticator για ση$\mu$αντικούς λογαριασ$\mu$ούς
  \item Πέντε καθη$\mu$ερινά $\mu$έτρα πρόληψης: κλείδω$\mu$α οθόνης+2FA, εύρεση κινητού, PIN SIM, όχι εφαρ$\mu$ογές αγνώστης προέλευσης, τακτικό backup
\end{itemize}

\end{tcbsummary}
